\newpage
\section{Missió 3: Pasturant el caos}\label{sec:mission3}

\begin{minipage}{0.45\linewidth}
A causa de diferències polítiques, i malgrat tots els esforços de l'assemblea de la colònia, cada base del planeta ha adoptat una tecnologia LAN (xarxa d'àrea local) diferent (!).
Tot i que les comunicacions dins de cada base són satisfactòries, ara per ara és impossible connectar directament dos ordinadors de dues bases diferents.
%
Un cop més, la colònia et demana ajuda.\\
\end{minipage}
\hspace{0.04\linewidth}
\begin{minipage}{0.5\linewidth}
\vspace{-1cm}\includegraphics[width=\linewidth]{colony_network_map.pdf}
\end{minipage}

\noindent
Gràcies als sacrificis de tothom, la colònia ja està en condicions de construir i instal·lar un nou router (encaminador) i dos busos cablejats.
Al diagrama següent es representen amb una línia més gruixuda.
%
% \noindent
El teu objectiu a llarg termini és fer possible la comunicació entre xarxes entre dos nodes qualssevol de dues bases qualssevol.
De moment, dissenyaràs la infraestructura física i lògica.\\

\subsection*{Detalls de l'informe}

Totes les bases fan servir el protocol IPv4, i hi ha diverses bases connectades tant al bus Nord com al bus Sud.
Cada base conté una única LAN. Coordina't amb els altres perquè \textit{no hi hagi dues LAN de la mateixa mida}.

\begin{enumerate}[itemsep=0.1cm]
\item Dins del teu grup:
\begin{itemize}[itemsep=0.1cm]
    \item Determina quantes adreces necessites perquè cada dispositiu es pugui identificar de manera única dins la teva LAN.
    Tria la mida de LAN mínima que satisfaci aquest requisit i quantifica l'``espai'' lliure disponible a
    la teva LAN.
    \item Indica quants dels 32~bits d'una adreça IPv4 necessites per identificar un sol dispositiu connectat a la LAN de la teva base.
    Quin concepte utilitzat a IP està més estretament relacionat amb aquest nombre? Per a què pot ser útil aquesta estructura?\\[0.2cm]
\end{itemize}

\begin{center}
\begin{bytefield}{18}
\bitbox{8}{Network ID} & \bitbox{15}{Host ID}\\
\end{bytefield}
\end{center}

\item En coordinació amb els altres grups de l'aula:
\begin{itemize}[itemsep=0.1cm]
\item Crea una LAN per al bus Nord i descriu totes les adreces físiques i lògiques que hi estan connectades.
\item Assegura't que els blocs d'adreces IPv4 assignats a les LAN siguin de mida diferent i no se solapin.
\item Fent servir la pissarra o un document compartit, fes una llista de totes les LAN, amb les seves mides i els rangs d'adreces (actualitzats).
\end{itemize}

\end{enumerate}

\subsection*{Detalls de la retroacció}

Segueix un procés de retroacció semblant al de les missions anteriors, i assegura't de tractar aquests punts:

\begin{enumerate}[itemsep=0.1cm]
\item Les LAN estan dimensionades correctament?
\item Hi ha cap col·lisió entre els blocs?
\item L'informe original extreia tota la informació sobre l'estructura de les adreces IP?
\end{enumerate}

